a. L'équation de la réaction est :

$$
2 \mathrm{H}_{2}(g)+\mathrm{O}_{2}(g) \longrightarrow 2 \mathrm{H}_{2} \mathrm{O}(\ell) \text {. }
$$

b. Les valeurs de $k$ sont :

$$
k_{\mathrm{H}_{2}}=\frac{n_{\mathrm{H}_{2}}^{i}}{2} \text { et } k_{\mathrm{O}_{2}}=\frac{n_{\mathrm{O}_{2}}^{i}}{1} .
$$

On noté $2 x$ ( $x$ : avancement) la quantité de matière d'eau formée. Les premières lignes sont complétées dans le tableau ci-dessous.

\begin{center}
\begin{tabular}{l|c|c|c|c}
\hline
\multicolumn{2}{c|}{Équation} & \multicolumn{2}{c}{$2 \mathrm{H}_{2}(g)+\mathrm{O}_{2}(g) \longrightarrow 2 \mathrm{H}_{2} \mathrm{O}(\ell)$} &  \\
\hline
\multirow{2}{*}{État initial (mol)} & Avancement & \multirow{2}{*}{$n_{\mathrm{H}_{2}}^{i}$} & $n_{\mathrm{O}_{2}}^{i}$ & 0 \\
\cline { 2 - 2 }
 & 0 &  &  &  \\
\hline
Au cours de la transformation (mol) & $x$ & $n_{\mathrm{H}_{2}}^{i}-2 x$ & $n_{\mathrm{O}_{2}}^{i}-x$ & $2 x$ \\
\hline
État final (mol) & $x_{\text {max }}=233$ & 0 & 78 & 466 \\
\hline
\end{tabular}
\end{center}

\begin{itemize}
  \item $n_{\mathrm{H}_{2}}^{i}=466 \mathrm{~mol} ; n_{\mathrm{O}_{2}}^{i}=311 \mathrm{~mol}$.
  \item $k_{\mathrm{H}_{2}}=\frac{466}{2}=233 ; k_{\mathrm{O}_{2}}=311$ : le réactif limitant est $\mathrm{H}_{2}$.
  \item Dans l'état final : $n_{\mathrm{H}_{2}}^{f}=0=466-2 x_{\max } \Rightarrow x_{\max }=233 \mathrm{~mol}$.
\end{itemize}

État final : $n_{\mathrm{H}_{2}}^{f}=0 ; n_{\mathrm{O}_{2}}^{f}=78 \mathrm{~mol} ; n_{\mathrm{H}_{2} \mathrm{O}}^{f}=466 \mathrm{~mol}$.


